\begin{frame}
    \frametitle{Ecuaciones de la estructura estelar}
    Describen la distribución de masa con simetría esférica
    \begin{align}
        \mathcolor{mygreen}{\diff{P\left(r\right)}{r}                                                                                     & =
            -\frac{Gm\left(r\right)}{r^{2}}
        \rho\left(r\right)}.\label{eq:hydrostatic}                                                                                            \\
        \mathcolor{myred}{
        \diff{m\left(r\right)}{r}                                                                                                         & =
            4\pi r^{2}\rho\left(r\right)}.\label{eq:mass}
        \shortintertext{Derivando
            $\mathcolor{myblue}{\dfrac{r^{2}}{\rho\left(r\right)}}\times\eqref{eq:hydrostatic}$
            respecto a $r$ da}
        \diff*{\left[
                \mathcolor{myblue}{\frac{r^{2}}{\rho\left(r\right)}}
                \mathcolor{mygreen}{\diff{P\left(r\right)}{r}}
        \right]}{r}                                                                                                                       & =
        -G
        \mathcolor{myred}{\diff{m\left(r\right)}{r}}.\label{eq:fusion}
        \shortintertext{Reemplazando~\eqref{eq:mass} en~\eqref{eq:fusion}}
        \mathcolor{myred}{\frac{1}{r^{2}}}
        \diff*{\left[\frac{r^{2}}{\rho\left(r\right)}\diff{P\left(r\right)}{r}\right]}{r}                                                 & =
        -\mathcolor{myred}{4\pi}G\mathcolor{myred}{\rho\left(r\right)}.\label{eq:total}
        \shortintertext{Un polítropo describe la relación en la que
            la presión depende de la densidad.}
        \mathcolor{mygreen}{
            P\left(r\right)=
            K{\left[\rho\left(r\right)\right]}^{\frac{n+1}{n}}\implies
        \diff{P\left(r\right)}{r}                                                                                                         & =
            \frac{\left(n+1\right)K}{n}
            {\left[\rho\left(r\right)\right]}^{\frac{1}{n}}
            \diff{\rho\left(r\right)}{r}
        }.\label{eq:polytrope}                                                                                                                \\
        \shortintertext{Reemplazando~\eqref{eq:polytrope} en~\eqref{eq:total}}
        \frac{1}{r^{2}}
        \frac{\mathcolor{mygreen}{\left(n+1\right)K}}{4\mathcolor{mygreen}{n}\pi G}
        \diff*{\left\{r^{2}\mathcolor{mygreen}{{\left[\rho\left(r\right)\right]}^{\frac{1-n}{n}}\diff{\rho\left(r\right)}{r}}\right\}}{r} & =
        -\rho\left(r\right).\label{eq:total2}
        \shortintertext{Introduzca la variable adimensional
            $\mathcolor{myblue}{r\left(\xi\right)\coloneqq a\xi}$,
            donde $\mathcolor{mygreen}{a^{2}=\frac{\left(n+1\right)K}{4\pi G}\lambda^{\frac{1-n}{n}}}$ y
            $\mathcolor{myred}{\rho\left(r\right)=\lambda{\left[\theta\left(r\right)\right]}^{n}}$.
            Reemplazando en~\eqref{eq:total2}}
        \mathcolor{myblue}{\frac{1}{a^{2}\xi^{2}}}
        \mathcolor{mygreen}{\frac{a^{2}\lambda^{\frac{n-1}{n}}}{n}}
        \diff*{\left\{
            \mathcolor{myblue}{a^{2}\xi^{2}}
            \mathcolor{myred}{
                \lambda^{\frac{1-n}{n}}
                {\left[\theta\left(\mathcolor{myblue}{a\xi}\right)\right]}^{1-n}
                \diff{\lambda{\left[\theta\left(\mathcolor{myblue}{a\xi}\right)\right]}^{n}}{\mathcolor{myblue}{a\xi}}
            }
        \right\}}{\mathcolor{myblue}{a\xi}}                                                                                               & =
        -\mathcolor{myred}{\lambda{\left[\theta\left(\mathcolor{myblue}{a\xi}\right)\right]}^{n}}.\notag
    \end{align}
\end{frame}

\begin{frame}
    Simplificando la última ecuación obtenemos
    \begin{equation}
        \boxed{
            \frac{1}{\xi^{2}}
            \diff*{\left(\xi^{2}\diff{\theta}{\xi}\right)}{\xi}=
            -\theta^{n}.\tag{Lane-Emden}
        }\label{eq:Lane-Emden}
    \end{equation}
    % Sea $k\geq 1$ un parámetro de forma
    \begin{block}{Problema de valor inicial singular de tipo Lane-Emden}
        \begin{equation}\label{eq:type-lane-emden}
            \begin{cases}\displaystyle
                \frac{1}{\xi^{m-1}}
                \diff*{\left(\xi^{m-1}\diff{\theta}{\xi}\right)}{\xi}=
                \left(g\circ\theta\right)\left(\xi\right)
                 & \text { para }0<\xi\leq R. \\[2mm]\displaystyle
                \theta\left(0\right)=
                \alpha,\quad
                \diff.|.{\theta}{\xi}[\xi=0]=
                0.
                 &
            \end{cases}
        \end{equation}
    \end{block}

    \vspace*{-2\baselineskip}\setlength\belowdisplayshortskip{0pt}
    \begin{equation}
        \boxed{
            \begin{aligned}
                \text{LHS} & :
                \frac{1}{\xi^{\alert{3}-1}}
                \diff*{\left(\xi^{\alert{3}-1}\diff{\theta}{\xi}\right)}{\xi}=
                \frac{1}{\xi^{2}}
                \left(
                2\xi\diff{\theta}{\xi}+
                \xi^{2}\diff[2]{\theta}{\xi}
                \right)=
                \frac{2}{\xi}
                \diff{\theta}{\xi}+
                \diff[2]{\theta}{\xi}. \\
                \text{RHS} & :
                \left(g\circ\theta\right)
                \left(\xi\right)=
                {\left[\theta\left(\xi\right)\right]}^{n},\quad
                n\geq 0.
            \end{aligned}
        }\tag{Lane-Emden}
    \end{equation}

    \begin{block}{Transformando~\eqref{eq:type-lane-emden} a un sistema EDO de primer orden}
        Considere el cambio de variable
        \begin{math}
            z\left(\xi\right)\leftarrow
            \diff.|.{\theta}{\xi}[\xi]=
            y\left(\xi\right)
        \end{math}.
    \end{block}

    \begin{equation}
        \begin{cases}\displaystyle
            \begin{bmatrix}
                \diff{y}{\xi} \\
                \diff{z}{\xi}
            \end{bmatrix}=
            \begin{bmatrix}
                z \\
                -\frac{2z}{\xi}-
                g\left(y\right)
            \end{bmatrix}
             & \text { para }0<\xi\leq R. \\[0.5cm]\displaystyle
            \begin{bmatrix}
                y\left(0\right) \\
                z\left(0\right)
            \end{bmatrix}=
            \begin{bmatrix}
                \alpha \\
                0
            \end{bmatrix}.
             &
        \end{cases}
    \end{equation}
\end{frame}